% ------------------------------------------------
\StartChapter{Introduction}{chapter:introduction}
% ------------------------------------------------

這是國立成功大學碩博士用畢業論文的LaTeX模版. 本模版是使用學校最新的畢業論文要求來設計(參考: 附錄 - 撰寫論文須知 P.\RefPage{appendix:thesis-spec}).

雖然本模版的目標是為了提供學生可以使用LaTeX來寫畢業論文. 但是各系所有各自的格式, 所以做了一個表列出已知的系所情況(參考: 附錄 - 可使用的系所 P.\RefPage{appendix:acceptable-dept}), 故請在使用前先留意自己的系所有沒有格式要求. 如果沒有, 則本模版應該是可以用來使用; 否則要看系所上的格式, 是否跟本模版有相同的寫法.

本模版分以下幾個主要部份來進行教學:

\begin{enumerate}
  \item 本模版使用教學: 下載, 檔案結構, 產生PDF的三種方式, 以及`conf/conf.tex'每一項怎麼填
  \item LaTeX編寫教學: 章節, 圖, 表, 公式, 符號表, 參考文獻和虛擬碼的寫法
  \item 附錄: 本專案保存的歷史學校文件, 本模版的起因(Chap.\RefTo{chapter:objective}), 以及``老師們的話''(Chap.\RefTo{chapter:words-from-professor})
\end{enumerate}

同學們只要閱讀完後, 把部份的檔案直接copy和修改內容, 應該很快就能上手本模版去寫自己的論文.

另外在附錄(Appendix)附上了一些本專案保存的歷史學校文件, 僅供參考. 學校及系所的現行規定永遠優先於本模版; 提交前請以學校系統及系所公佈的最新版本為準.

% ------------------------------------------------
\newpage

\begin{description}
  \item[版權 License]\hfill\\
  詳細請看`LICENSE'這檔案中的條款說明.\\

  \InsertFigure
    [scale=0.8,
      caption={CC Attribution-NonCommercial-ShareAlike License},
      label={fig:appendix:by-nc}]
    {./example/introduction/pic/by-nc-sa.png}

    本著作(ncku-thesis-template-latex\RefBib{web:this-project:github})採用創用 CC 姓名標示-非商業性-相同方式分享 4.0 授權條款.

    This work(ncku-thesis-template-latex\RefBib{web:this-project:github}) is licensed under Creative Commons Attribution-NonCommercial-ShareAlike 4.0 International License.\\

    這授權只涵蓋本模版本身的檔案. 同學用本模版撰寫的論文內容, 圖表, 資料和產生出來的PDF都屬於同學自己, 本專案不對其主張任何權利, 也不要求論文採用相同的授權.\\

  而本模版所使用到的國立成功大學浮水印則由國立成功大學擁有\textbf{所有}相關的權利. 故如使用這浮水印到論文以外的應用, 請跟`成大圖書館 系統管理組-數位論文小組'聯絡.\\

  \item[版本修改 ChangeLog]\hfill\\
  詳細請看`ChangeLog.md'這檔案中的說明.
\end{description}

% ------------------------------------------------
\EndChapter
% ------------------------------------------------
